\uuid{YsX7}
\titre{Symétrie de la loi normale centrée réduite}
\auteur{Maxime Nguyen}
\organisation{AMSCC}
\datecreate{2026-11-24}
\niveau{L2}
\module{Probabilité}
\chapitre{Probabilité continue}
\sousChapitre{Loi normale}
\difficulte{1}
\theme{loi normale centrée réduite, symétrie, fonction de répartition}

\contenu{
\begin{question}
Soit $Z$ une variable de loi normale centrée réduite $\mathcal{N}(0,1)$, $\Phi$ sa fonction de répartition ($\Phi(t) = \mathbb{P}(Z \leq t)$) et $t > 0$. Quelles égalités sont vraies ?
\begin{answers}
\bad{$\Phi(-t) = -\Phi(t)$}
\good{$\Phi(-t) = 1 - \Phi(t)$}
\bad{$\Phi(0) = 0$}
\good{$\mathbb{P}(|Z| \leq t) = 2\Phi(t) - 1$}
\good{$\mathbb{P}(Z \geq t) = \Phi(-t)$}
\end{answers}
\begin{explanations}
A donnerait une probabilité négative. C : $Z$
est centrée et sa densité est paire, donc $\Phi(0) = \mathbb{P}(Z \leq 0) = \frac12$ ;
c'est la densité, pas $\Phi$, qui est symétrique par rapport à $0$. Par
symétrie de la densité, $\mathbb{P}(Z \leq -t) = \mathbb{P}(Z \geq t) = 1 - \Phi(t)$,
et $\mathbb{P}(-t \leq Z \leq t) = \Phi(t) - \Phi(-t) = 2\Phi(t) - 1$.

Contre-sens visé : confondre la symétrie de la densité (paire) avec une
imparité de $\Phi$ (A, C). Le réflexe : faire un dessin de la cloche et
colorier l'aire.
\end{explanations}
\end{question}
}
